% ============================================================================
% Related work v2.4 — replacement block for main.tex (\section{Related work}).
% Synced from the author's edited related-work-v2-zh.md.
% v2.4 changes: (i) section is now THREE paragraphs, mirroring the author's
% edit — the subgraph-sampling paragraph was removed on top of the earlier
% removals (budgeted acquisition, complexity); keys graphsage/crawl moved to
% UNUSED below; (ii) the "learned sequential readers" paragraph profiles the
% four closest works (GAM, RWE-SGNN, AgentNet, RAW) by behavior, task, and
% protocol violations, as fact-checked against the source full texts;
% (iii) no forward references: the T0/T1/T2 tiers are defined in the
% introduction, so invoking them here is a backward reference.
%
% Conventions kept from main.tex:
%   \refadd{key}     citation placeholder (bibliography intentionally omitted)
% ============================================================================

\section{Related work}
\label{sec:related}

\paragraph{Graph classification and pooling.}
Message-passing classifiers such as GIN consume the whole graph \refadd{gin}, and pooling methods sparsify this computation by scoring nodes with task-aware attention on the full graph, from differentiable cluster assignment to self-attention top-$k$ selection \refadd{diffpool,sagpool,pooling-survey}. All of these answer \emph{which nodes matter} once the full graph has been read; under a node budget the same question must be answered \emph{before} the graph is visible, from the frontier alone.

\paragraph{Graph explainability.}
Explanation methods identify rationales for a trained model's prediction on an already-observed graph, from the edge masks of GNNExplainer to the connected rationale subgraphs of SubgraphX \refadd{gnnexplainer,subgraphx}. Our setting inverts the direction of this step: evidence must be acquired \emph{before} there is a prediction to explain, so explanation-style scores are unavailable by construction under partial observation.

\paragraph{Learned sequential readers.}
Closest to our setting are agents that read a graph selectively under a budget. \textbf{GAM} trains, by reinforcement learning, a recurrent walker for budgeted whole-graph classification: a rank network inspects the current node's neighbors and samples the next move, an LSTM integrates the reading history, and multiple walkers can pool their histories into shared memory \refadd{gam}. \textbf{RWE-SGNN} keeps the walk but learns it by DQN, rewarding the induced drop of a downstream loss; its separated sampler and predictor generate class-relevant substructures for self-supervised graph classification \refadd{rwesgnn}. \textbf{AgentNet} replaces reinforcement learning with straight-through Gumbel transitions: a population of identification-carrying agents walks the graph differentiably and sublinearly, exchanging state with the visited node at every step \refadd{agentnet}. \textbf{RAW} recasts the walk as a POMDP for semi-supervised \emph{node} classification, admitting edge content into the observation and aggregating the features of neighbors whose relevance passes a threshold \refadd{raw}. The four tasks bracket ours: RAW shields a single node's prediction while the rest of the graph is treated as observable, whereas the other three share our graph-level objective. All of them, however, violate the discipline our protocol enforces. First, unrevealed content enters the decision path and even the predictive representation: GAM inspects neighbors' content before moving to them, RAW aggregates never-visited neighbors' features, and AgentNet absorbs neighborhood states at every step while its differentiable transitions backpropagate into unvisited candidates' embeddings. Second, none maintains a single-seed, monotonically growing revealed set: budgets are counted in walk steps or agents, GAM may revisit nodes, RWE-SGNN launches multiple walks from graph-wide top-scoring entry points, and AgentNet scatters its agents across the whole graph. Third, GAM, RAW, and AgentNet read the prediction from the same recurrent state that steers selection; RWE-SGNN's sampler--predictor split travels the same road as our decoupling, yet carries no reveal discipline. Our reader, by construction, guarantees that the predictive path is a function of the revealed set alone.

% ----------------------------------------------------------------------------
% Citation keys -> full references (for the future bibliography).
%
% USED in this section (10 keys):
%   gin         Xu, Hu, Leskovec, Jegelka, "How Powerful are Graph Neural Networks?" ICLR 2019
%   diffpool    Ying et al., "Hierarchical Graph Representation Learning with Differentiable Pooling," NeurIPS 2018
%   sagpool     Lee, Lee, Yoo, Choi, Kahng, "Self-Attention Graph Pooling," ICML 2019
%   pooling-survey   Grattarola, Zambon, Bianchi, Alippi, "Understanding Pooling in Graph Neural Networks," IEEE TNNLS 2022
%   gnnexplainer Ying et al., "GNNExplainer," NeurIPS 2019
%   subgraphx   Yuan, Yu, Gui, Ji, "On Explainability of GNNs via Subgraph Explorations," ICML 2021
%   gam         Lee et al., "Graph Classification using Structural Attention," KDD 2018
%   rwesgnn     "Self-supervised Subgraph Neural Network with Deep Reinforcement Walk Exploration," arXiv 2502.01809, 2025
%   agentnet    Martinkus et al., "Agent-based Graph Neural Networks," ICLR 2023
%   raw         Akujuobi et al., "Recurrent Attention Walk for Semi-supervised Classification," WSDM 2020
%
% UNUSED here (removed paragraphs; reuse if the material moves elsewhere):
%   graphsage   Hamilton, Ying, Leskovec, "Inductive Representation Learning on Large Graphs," NeurIPS 2017
%   crawl       Tönshoff, Ritzert, Rosenkranz, Grohe, "Graph Learning with 1D Convolutions on Random Walks," AISTATS 2021
%   gpa         Hu et al., "Graph Policy Network for Transferable Active Learning on Graphs," NeurIPS 2020
%   igp         Zhang et al., "Information Gain Propagation: A New Way to Graph Active Learning with Soft Labels," ICLR 2022
%   gflowgnn    Li et al., "Generative Flow Networks for Precise Reward-Oriented Active Learning on Graphs," arXiv 2304.11989, 2023
%   afa-shim    Shim, Hwang, Yang, "Joint Active Feature Acquisition and Classification with Variable-Size Set Encoding," NeurIPS 2018
%   ram         Mnih, Larochelle, Hinton, "Recurrent Models of Visual Attention," NeurIPS 2014
%   grl-survey  Nie et al., "Reinforcement Learning on Graphs: A Survey," IEEE TETCI 2022
%   games-against-nature  Papadimitriou, "Games against Nature," JCSS 1985
%   pomdp-complexity      Papadimitriou & Tsitsiklis, "The Complexity of Markov Decision Processes," Mathematics of OR 1987
%   ssat        Majercik & Littman, "Stochastic Boolean Satisfiability," AAAI 1998
%   adaptive-submodularity  Golovin & Krause, "Adaptive Submodularity: Theory and Applications in Active Learning and Stochastic Optimization," JAIR 2011
% ----------------------------------------------------------------------------
